\documentclass[11pt]{article}
\usepackage[a4paper,margin=1in]{geometry}
\usepackage{amsmath,amssymb,mathtools,bm}
\usepackage{enumitem,booktabs,array}
\usepackage{tcolorbox}
\usepackage{hyperref}
\usepackage[protrusion=true,expansion=false]{microtype}
\hypersetup{colorlinks=true,linkcolor=blue,urlcolor=blue,citecolor=blue}
\setlist[itemize]{topsep=2pt,itemsep=2pt}
\setlist[enumerate]{topsep=2pt,itemsep=2pt}
\newtcolorbox{keybox}{colback=blue!3!white,colframe=blue!40!black,boxrule=0.4pt,arc=1mm}
\newtcolorbox{alertbox}{colback=red!3!white,colframe=red!40!black,boxrule=0.4pt,arc=1mm}
\newtcolorbox{answerbox}{colback=green!3!white,colframe=green!40!black,boxrule=0.4pt,arc=1mm}
\newtcolorbox{drillbox}{colback=yellow!5!white,colframe=orange!60!black,boxrule=0.4pt,arc=1mm}
\newcommand{\EV}{\mathbb{E}}
\newcommand{\Var}{\mathrm{Var}}
\newcommand{\Cov}{\mathrm{Cov}}
\newcommand{\blank}[1]{\underline{\hspace{#1}}}

\title{W12-03: Bernstein, Gustafson \& Lewis (2019, JFE) \\
\large ``Disaster on the Horizon: The Price Effect of Sea Level Rise'' --- Active Recall Workbook}
\author{Banking \& Risk Management --- Exam Prep}
\date{\today}
\begin{document}\maketitle

\begin{keybox}
\textbf{Paper in one sentence.} US coastal homes exposed to sea-level-rise (SLR) risk (projected inundation under 6 ft SLR scenario) trade at a $\sim$7\% discount relative to otherwise-identical non-exposed homes within the same ZIP, and the discount is concentrated in areas with higher climate-concern populations --- housing markets \emph{do} price SLR risk, at least in part.

\textbf{Placement.} Headline climate-pricing paper in real estate. Contrasts with Murfin-Spiegel (W13-01), which finds a much smaller effect after controlling for subsidence. Companion paper Bernstein-Billings-Gustafson-Lewis (W12-04) explores \emph{whose} beliefs drive the discount.
\end{keybox}

\section*{Round 1: Complete Walk-Through}

\subsection*{1.1 Exposure and data}
NOAA SLR inundation maps (6 ft scenario) define exposure at the parcel level. Sample: 460{,}000 coastal home sales, 2007--2016, across US Atlantic and Gulf coasts.

\subsection*{1.2 Identification}
Hedonic regression within ZIP-by-quarter cells:
\[
\ln P_i=\alpha_{z,t}+\beta\cdot\text{SLRexposed}_i+X_i'\gamma+\varepsilon_i,
\]
with $X_i$ including physical distance to coast, elevation, and standard hedonic characteristics (square footage, age, bedrooms, etc.). Identification comes from comparing SLR-inundation-exposed vs.\ non-exposed homes \emph{within} the same ZIP-quarter, holding ocean amenity fixed.

\subsection*{1.3 Headline results}
\begin{itemize}
\item $\beta\approx -7\%$: SLR-exposed homes trade at a $\sim$7\% discount.
\item Discount emerges in 2007 and widens mildly over time.
\item Stronger in ZIPs with higher climate-concern populations (proxied by Google Trends and NOAA surveys).
\item Effect absent in matched non-coastal samples (placebo).
\end{itemize}

\subsection*{1.4 Mechanism}
\begin{itemize}
\item Buyers internalise long-horizon SLR cost.
\item Sophisticated investors/second-home owners drive the discount; primary owner-occupiers less so.
\item Climate-belief channel: places where residents believe in climate change discount more.
\end{itemize}

\subsection*{1.5 Caveats}
\begin{alertbox}
\begin{itemize}
\item Amenity value of coastline is strong; separating SLR risk from location is challenging.
\item Ground subsidence may correlate with SLR exposure, confounding the estimate (see Murfin-Spiegel).
\item 7\% is small relative to present-value of projected SLR damage over long horizons (suggests under-pricing).
\end{itemize}
\end{alertbox}

\section*{Round 2: Level 1 Fill-in}
\begin{enumerate}
\item Exposure: NOAA \blank{1cm}-ft SLR inundation projection.
\item Sample: $\sim$\blank{3cm} coastal home sales, 2007--2016.
\item Discount: $\sim$\blank{1cm}\%.
\item Stronger in areas with higher \blank{1.5cm}.
\item Identification: within \blank{1cm}-\blank{1cm} FE.
\end{enumerate}

\section*{Round 3: Level 2 Prompts}
\begin{enumerate}
\item Write the hedonic spec and state the identifying assumption.
\item Why does within-ZIP-quarter identification control for amenity?
\item Interpret the climate-concern heterogeneity.
\item Contrast BGL's $-7\%$ with Murfin-Spiegel's smaller estimate.
\end{enumerate}

\section*{Round 4: Minimal Scaffolding}
From a blank page: (i) exposure, (ii) hedonic spec, (iii) magnitude, (iv) heterogeneity, (v) caveats.

\section*{Round 5: Blank Page}
Essay: is climate risk priced in coastal real estate? BGL evidence.

\vspace{6pt}
\begin{drillbox}
\textbf{Speed Drill.} (1) Exposure measure. (2) Sample size/years. (3) Discount magnitude. (4) FE structure. (5) Heterogeneity source.
\end{drillbox}

\section*{Quick Reference \& Answer Key}
\begin{answerbox}
\begin{itemize}
\item \textbf{Exposure.} NOAA 6 ft SLR inundation.
\item \textbf{Discount.} $\sim$7\% within ZIP-quarter.
\item \textbf{Sample.} 460{,}000 coastal homes, 2007--2016.
\item \textbf{Heterogeneity.} Stronger in climate-concerned ZIPs.
\end{itemize}
\end{answerbox}

\begin{keybox}
\textbf{30-second exam pitch.} Bernstein, Gustafson \& Lewis (2019) show that sea-level-rise risk is priced in US coastal real estate. Using NOAA 6-foot SLR inundation maps to define parcel-level exposure and within-ZIP-quarter hedonic regressions on 460{,}000 coastal home sales 2007--2016, they estimate that SLR-exposed homes sell at a $\sim$7\% discount relative to otherwise-identical non-exposed homes. The discount is stronger in areas with higher climate concern (proxied by Google Trends and survey data). Mechanism: buyers internalise long-run SLR damages; sophisticated investors are most sensitive. The result is the canonical climate-pricing finding in housing but is debated --- Murfin-Spiegel (W13-01) get a much smaller estimate after controlling for subsidence, and the follow-on BBGL (W12-04) shows beliefs, not just exposure, drive the discount.
\end{keybox}

\end{document}
